\chapter{A circulação de capital e a mercadoria-energia}

\setlength{\epigraphwidth}{0.7\textwidth}

\epigraph{Vastas redes de trilhos brilhantes, estendendo-se por aterros, pontes e viadutos, cortes e túneis de até dezesseis quilômetros de extensão, e cruzando passagens montanhosas tão elevadas quanto os principais picos dos Alpes, as ferrovias constituíram, em seu conjunto, o maior empreendimento de obras públicas jamais realizado pelo homem. Elas empregavam mais trabalhadores do que qualquer outro tipo de atividade industrial. Chegavam tanto aos centros das grandes cidades, onde suas conquistas triunfais eram celebradas em estações ferroviárias igualmente grandiosas e imponentes, quanto às áreas rurais mais remotas, onde nenhum outro vestígio da civilização do século XIX conseguia penetrar.}{\textcite[p. 27]{Hobsbawm1987}}

\epigraph{Nós cantaremos as grandes multidões agitadas pelo trabalho, pelo prazer ou pela sublevação; cantaremos as marés multicoloridas e polifônicas das revoluções nas capitais modernas; cantaremos o vibrante fervor noturno dos arsenais e dos estaleiros incendiados por violentas luas elétricas; as estações esfomeadas, devoradoras de serpentes esfumaçantes; as oficinas penduradas às nuvens pelos fios contorcidos de suas fumaças; as pontes, semelhantes a ginastas gigantes que cavalgam os rios, faiscantes ao sol com um luzir de facas; os navios a vapor aventurosos que farejam o horizonte, as locomotivas de amplo seio [\textit{dall'ampio petto}], que cavalgam sobre os trilhos como enormes cavalos de aço freados por tubos; e o voo rasante dos aviões, cuja hélice freme ao vento como uma bandeira e parece aplaudir como uma multidão entusiasta.}{\textcite{Marinetti1909}}

No que diz respeito ao efeito \textit{rebound/backfire}, a economia neoclássica busca a causa, ou uma das causas mais importantes, do aumento do consumo energético no aumento da eficiência de serviços energéticos. O que exatamente tais ganhos querem dizer, na prática, tende a ser algo ofuscado, tal que seu foco na esfera da circulação se torna justificado: pode-se falar livremente da eficiência \textit{dos serviços energéticos usufruídos}. 

Tal foco nestes ganhos de eficiência ao nível da circulação acaba por abstrair o papel crucial que o aumento das forças produtivas --- leia-se: na esfera da produção --- desempenham sobre o nível de consumo da mercadoria-energia. De fato, mistifica-se a diferença entre seu usufruto produtivo --- na produção de mercadorias, produção de valor --- e seu uso improdutivo.


A tendência ao aumento das forças produtivas, no modo de produção capitalista, carrega consigo uma contradição: enquanto permite uma produção aumentada de mercadorias, exerce uma pressão maior sobre o mercado; ao mesmo tempo que facilita sua produção, dificulta sua incorporação na circulação e na esfera do consumo.\footnote{``Crescem, com os avanços no processo produtivo, as exigências sobre o consumo e a necessidade de ampliar a escala do consumo não só eventualmente, mas sistematicamente.'' \parencite[p. 156]{SaBarreto2018}}

Embora o capital seja, em essência, valor que se valoriza, trabalho abstrato que busca se expandir, ele ainda se assenta sobre valores de uso, finitos por natureza e, portanto, contraditoriamente, possíveis obstáculos ao movimento do capital. Como resumido por Sá Barreto,
\begin{quote}
    O valor de uso impõe, portanto, pela sua própria natureza, um dado limite à realização do mais-valor produzido, limite que é determinado pelas necessidades às quais atende e pelo tempo durante o qual pode atendê-las sem que seja necessária a sua substituição. Por isso, no curso do desenvolvimento do sistema, que é orientado para a expansão do valor e apenas indiretamente para a satisfação de necessidades, deve também o valor de uso assumir formas que se adaptem a este objetivo primordial. Deve, por conseguinte, o valor de uso crescentemente assumir formas evanescentes, fugazes. \parencite[p. 157]{SaBarreto2018}
\end{quote}

\section{Capital de comércio}
É do interesse do capital industrial sempre aumentar suas forças produtivas, diminuindo seu tempo de rotação, em particular através da diminuição do tempo de produção. Ele não consegue diminuir o tempo de circulação por si próprio, ou ao menos não é de seu interesse imediato separar recursos para dispêndio em \textit{faux frais}. 

A fagocitose da esfera da circulação pelo próprio capital requer que ele se apodere deste âmbito social, que se torne senhor daquilo que lhe aparecia --- enquanto capital industrial, i.e. neste nível de abstração (?) --- como dado, como força da natureza (??). É dessa forma que vem a ser sua mitose em algo qualitativamente diferente: o \textit{capital comercial}, a circulação de capital enquanto algo sob o domínio dele próprio.

Obviamente, o processo de redução do tempo de circulação do capital, sob demanda do capital industrial, não poderia ser resultado do altruísmo dos caixeiros-viajantes, não só porque eles agem por seu ``amor próprio'' como por ser algo contingente ao capital. Dessa forma, ele precisa introjetar-se ali, impor seu ímpeto de dentro para fora nesta esfera --- em suma, reconstruí-la à sua imagem e semelhança. Com isso, a circulação de capital passa a agir \textit{também enquanto capital}; passa a ser, agora, capital que faz capital circular.\footnote{Não se trata de ``capital circulante'', porém! Provavelmente é por isso que Marx usou, certas vezes, o termo ``capital líquido'' para falar da categoria de capital circulante.}

O processo de circulação, porém, não produz valor \textit{per se}, meramente \textit{circula} o valor já produzido e existente. Seus agentes, porém, não se importam com essas tecnicalidades: eles despenderam cérebros, nervos, músculos etc. em seu \textit{métier}, e demandam uma remuneração como recompensa por seu trabalho. 

Na esfera da produção de mercadorias, desdobra-se uma contradição entre o valor que o capital \textit{produz} e o de que ele se \textit{apropria}. Na esfera da circulação de mercadorias, tal contradição desdobra-se ainda mais: nos setores que não são produtores de mercadorias, mas que são meramente auxiliares destes setores, o dispêndio de trabalho, o \textit{lucro comercial} etc. são ``custos improdutivos'' do ponto de vista do capital industrial, embora apareça --- e isso não é uma mera mistificação --- como o pagamento devido pelo ``suor de sua fronte''\footnote{Gên. 3:19.}. 